\documentclass[11pt,a4paper]{article}

\usepackage[T1]{fontenc}
\usepackage{lmodern}
\usepackage{amsmath,amssymb}
\usepackage[margin=2.5cm]{geometry}
\usepackage{booktabs}
\usepackage[colorlinks=true,linkcolor=blue,urlcolor=blue]{hyperref}

\newcommand{\Mqq}{M_{q\bar q}^2}
\newcommand{\dd}{\mathrm{d}}
\newcommand{\xP}{x_{\mathbb{P}}}

\title{The finite-energy 5D cross section\\integrated by parts in $\Mqq$}
\author{}
\date{\today}

\begin{document}
\maketitle

\begin{abstract}
  The finite-energy (FE) constrained DIS cross section with the dipole evolved
  to the rapidity of $\xP$ is a five-dimensional integral whose $\Mqq$ integrand
  oscillates over the whole kinematic range. Integrating by parts in $\Mqq$
  rewrites it exactly as the four-dimensional FE cross section with the dipole
  at $\xP = 1$, plus a correction proportional to the rapidity derivative of
  the target amplitude. The correction vanishes wherever the dipole does not
  depend on rapidity, which for a BK solution starting at $x_0$ is most of the
  $\Mqq$ range. Numerically, this reduces the Monte Carlo error at a fixed
  number of points by a factor 12--17, and it makes adaptive cubature
  practical.
\end{abstract}

\section{The cross section}

We follow the implementation in
\texttt{BK\_FE\_5D\_evolve\_in\_xP.py}. For one quark flavour of mass $m$ and
polarization $\lambda = L, T$,
\begin{equation}
  \sigma_\lambda^{\text{5D}}
  = \int \dd u \, \dd u' \, \dd\theta \, \dd z
    \int_{4m^2}^{W^2} \dd\Mqq \;
    P_\lambda(u, u', \theta, z) \;
    T\big(Y(\Mqq)\big) \;
    \Theta\!\left(\Mqq - \frac{m^2}{z(1-z)}\right)
    \frac{z(1-z)}{4\pi} \, J_0(k r) ,
  \label{eq:sigma5d}
\end{equation}
where $u$ and $u'$ are the dipole sizes in the amplitude and its conjugate,
$\theta$ is the angle between them, and
\begin{align}
  r^2 &= |\mathbf{u} - \mathbf{u}'|^2 = u^2 + u'^2 - 2 u u' \cos\theta , &
  k^2 &= z(1-z) \, \Mqq - m^2 , \\
  W^2 &= Q^2 \left(\frac{1}{x_B} - 1\right) , &
  Y(\Mqq) &= \ln \frac{x_0}{\xP} , \qquad
  \xP = \frac{\Mqq + Q^2}{W^2 + Q^2} .
\end{align}
The prefactor collects everything that does not depend on the target:
\begin{equation}
  P_\lambda(u, u', \theta, z)
  = \frac{\sigma_0}{2} \, \frac{1}{4\pi} \,
    \frac{2\pi \, u u'}{z(1-z)} \,
    \big|\psi_\lambda(u, u', \theta, z)\big|^2 ,
\end{equation}
with $|\psi_\lambda|^2$ the product of the light-cone wavefunction and its
conjugate. The target amplitude
\begin{equation}
  T(Y) = 1 - S_Y(u) - S_Y(u') + S_{4,Y}(u, u', \theta, z)
\end{equation}
is built from the dipole $S_Y$ and the quadrupole $S_{4,Y}$, all at rapidity
$Y$.

The structure that matters below is that \textbf{only $T$ depends on $\Mqq$},
through the rapidity $Y(\Mqq)$, and that the phase-space kernel is $J_0(kr)$
with $k^2$ linear in $\Mqq$.

For comparison, the 4D cross section with the dipole evolved to $x_B$ is
\begin{equation}
  \sigma_\lambda^{\text{4D}}[Y]
  = \int \dd u \, \dd u' \, \dd\theta \, \dd z \;
    P_\lambda(u, u', \theta, z) \; T(Y) \;
    \frac{\zeta J_1(\zeta)}{2\pi r^2} ,
  \qquad
  \zeta = r \sqrt{z(1-z) \, M_{\max}^2 - m^2} ,
  \label{eq:sigma4d}
\end{equation}
with $M_{\max}^2 = Q^2 (1 - x_B)/x_B = W^2$ and the integrand set to zero where
the square root is not real. \texttt{BK\_FE\_4D\_evolve\_in\_xB.py} uses
$Y = \ln(x_0/x_B)$; here the rapidity is left as an argument.

\section{Integration by parts}

Fix $u$, $u'$, $\theta$ and $z$. The $\Theta$ function starts the $\Mqq$
integral at the threshold $\Mqq = m^2/[z(1-z)] \geq 4m^2$, where $k = 0$. The
upper limit $\Mqq = W^2$ corresponds to $k = K$, with
\begin{equation}
  K^2 = z(1-z) \, W^2 - m^2 .
\end{equation}
If $K^2 \leq 0$ the integrand vanishes, and there is nothing to do. Otherwise,
change variables to $s = k^2$, so that $\dd s = z(1-z) \, \dd\Mqq$. The
$\Mqq$ integral in \eqref{eq:sigma5d} becomes
\begin{equation}
  I \equiv \int_{4m^2}^{W^2} \dd\Mqq \;
    \Theta(\ldots) \, \frac{z(1-z)}{4\pi} \, J_0(kr) \, T
  = \frac{1}{4\pi} \int_0^{K^2} \dd s \; J_0\big(\sqrt{s}\, r\big) \, T .
\end{equation}

The kernel has an elementary antiderivative in $s$. From
$\frac{\dd}{\dd x} \big[ x J_1(x) \big] = x J_0(x)$ it follows that
\begin{equation}
  F(s) \equiv \frac{2\sqrt{s}\, J_1(\sqrt{s}\, r)}{r} ,
  \qquad
  \frac{\dd F}{\dd s} = J_0\big(\sqrt{s}\, r\big) ,
  \qquad
  F(0) = 0 .
\end{equation}
Integrating by parts,
\begin{equation}
  I = \frac{1}{4\pi} \Big[ F(s) \, T \Big]_{0}^{K^2}
    - \frac{1}{4\pi} \int_0^{K^2} \dd s \; F(s) \, \frac{\dd T}{\dd s} .
\end{equation}
The lower boundary term vanishes because $F(0) = 0$. In the remaining integral,
$\dd s \, (\dd T / \dd s) = \dd\Mqq \, (\dd T / \dd \Mqq)$, and by the chain rule
\begin{equation}
  \frac{\dd T}{\dd \Mqq}
  = \frac{\dd T}{\dd Y} \, \frac{\dd Y}{\dd \Mqq}
  = - \frac{1}{\Mqq + Q^2} \, \frac{\dd T}{\dd Y} .
\end{equation}
Therefore
\begin{equation}
  \boxed{
  I = \frac{K J_1(K r)}{2\pi r} \, T(Y_W)
    + \frac{1}{4\pi} \int \dd\Mqq \;
      \frac{2 k J_1(k r)}{r} \,
      \frac{1}{\Mqq + Q^2} \,
      \frac{\dd T}{\dd Y}\bigg|_{Y(\Mqq)}
  }
  \label{eq:ibp}
\end{equation}
where the integral runs over the same range as before, from threshold to $W^2$,
and $Y_W = Y(W^2) = \ln x_0$, the rapidity at $\xP = 1$.

\section{Result}

\paragraph{The boundary term is the 4D cross section.}
With $\zeta = K r$, the first term of \eqref{eq:ibp} is
\begin{equation}
  \frac{K J_1(K r)}{2\pi r} = \frac{\zeta J_1(\zeta)}{2\pi r^2} ,
  \qquad
  \zeta = r \sqrt{z(1-z) \, W^2 - m^2} ,
\end{equation}
which is exactly the kernel of \eqref{eq:sigma4d}, because $M_{\max}^2 = W^2$.
Multiplying by $P_\lambda$ and integrating over $u$, $u'$, $\theta$ and $z$
gives
\begin{equation}
  \boxed{
  \sigma_\lambda^{\text{5D}}
  = \sigma_\lambda^{\text{4D}}\big[Y = \ln x_0\big]
  + \Delta\sigma_\lambda
  }
\end{equation}
with the correction
\begin{equation}
  \begin{split}
  \Delta\sigma_\lambda
  = \int \dd u \, \dd u' \, \dd\theta \, \dd z \, \dd\Mqq \;
    &P_\lambda(u, u', \theta, z) \;
    \Theta\!\left(\Mqq - \frac{m^2}{z(1-z)}\right) \\
    &\times \frac{1}{4\pi} \, \frac{2 k J_1(k r)}{r} \,
    \frac{1}{\Mqq + Q^2} \,
    \frac{\dd T}{\dd Y}\bigg|_{Y(\Mqq)} .
  \end{split}
  \label{eq:correction}
\end{equation}

\paragraph{Interpretation.}
The 5D cross section equals the 4D one evaluated with the target at $\xP = 1$,
plus a term driven by how fast the evolution changes the target amplitude,
$\partial T / \partial Y$, weighted by a smooth kernel. At $\xP = 1$ we have
$Y = \ln x_0 < 0$. A BK solution starting at $x_0$ is not evolved backwards,
and the code evaluates the dipole at $Y < 0$ at its initial condition $Y = 0$.
So the boundary term is the 4D cross section with the unevolved dipole.

\paragraph{The correction lives at small $\Mqq$.}
The derivative $\dd T / \dd Y$ vanishes wherever the dipole does not depend on
$Y$. For a BK solution tabulated in $Y \in [Y_{\min}, Y_{\max}]$, with the
dipole frozen outside, the integrand of \eqref{eq:correction} is non-zero only
for
\begin{equation}
  M^2(Y_{\max}) \leq \Mqq \leq M^2(Y_{\min}) ,
  \qquad
  M^2(Y) = x_0 \left(W^2 + Q^2\right) e^{-Y} - Q^2 .
\end{equation}
With $Y_{\min} = 0$, the upper end corresponds to $\xP = x_0$. For
$Q^2 = 2\ \text{GeV}^2$, $x_B = 10^{-3}$ and $x_0 = 0.01$, the correction
covers $\Mqq \lesssim 18\ \text{GeV}^2$, while the original integral extends
to $W^2 \approx 2000\ \text{GeV}^2$.

\paragraph{Validity.}
Integration by parts requires $T$ to be absolutely continuous in $\Mqq$, not
smooth. Freezing the dipole at the edges of its rapidity range is continuous;
$\dd T / \dd Y$ then jumps there, which is harmless. The kernel of the
correction is regular at $r \to 0$, where $2kJ_1(kr)/r \to k^2$.

The derivation assumes that $T$ is the only $\Mqq$-dependent factor and that
the kernel is $J_0(kr)$ with $k^2$ linear in $\Mqq$. If either changes, for
example through an $\Mqq$-dependent factor elsewhere, the derivation has to be
redone.

\section{Why this helps numerically}

At fixed $u$, $u'$, $\theta$, $z$, the original integrand oscillates in $\Mqq$
through $J_0(kr)$ over the whole range up to $W^2$. Integrated over $\Mqq$,
these oscillations produce the kernel $\zeta J_1(\zeta)/(2\pi r^2)$ of the 4D
cross section. That kernel is sharply peaked at $r = 0$, with a width of about
$1/\sqrt{z(1-z)W^2}$, roughly $0.05\ \text{GeV}^{-1}$ in the example above. The
peak emerges from a near-cancellation between large positive and negative
contributions. A Monte Carlo estimate samples both and has a correspondingly
large variance.

After the integration by parts, that cancellation is done analytically, in the
antiderivative $F(s)$. What is left to integrate numerically is
\begin{itemize}
  \item the 4D boundary term. Its peak at $r = 0$ is handled by integrating over
        $\mathbf{u}$ and the relative vector $\mathbf{r} = \mathbf{u} -
        \mathbf{u}'$ in polar coordinates, with $|\mathbf{r}|$ on a logarithmic
        scale, instead of over $u$, $u'$ and $\theta$;
  \item the 5D correction, over a short $\Mqq$ range where $k$ is small and the
        kernel barely oscillates.
\end{itemize}

Table~\ref{tab:numbers} shows the effect for $Q^2 = 2\ \text{GeV}^2$,
$x_B = 10^{-3}$, $m = 0.14\ \text{GeV}$, $Z_f^2 = 2/3$, the MV-initial-condition
BK solution of \texttt{data/mv.dat}, and the finite-$N_c$ quadrupole. The
numbers are structure functions $F_\lambda = Q^2 \sigma_\lambda/(4\pi^2
\alpha_{\text{em}})$.
\begin{table}[h]
  \centering
  \small
  \begin{tabular}{lll}
    \toprule
    & $F_L$ & $F_T$ \\
    \midrule
    Vegas, $10^6$ points, original form          & $0.0752 \pm 0.0020$ & $0.2712 \pm 0.0044$ \\
    Vegas, $10^6$ points, integrated by parts    & $0.07620 \pm 0.00012$ & $0.27574 \pm 0.00038$ \\
    \quad of which the boundary term            & $0.06184 \pm 0.00012$ & $0.24603 \pm 0.00038$ \\
    \quad of which the correction               & $0.014362 \pm 0.000006$ & $0.029712 \pm 0.000011$ \\
    Cuhre, $\epsilon_{\text{rel}} = 10^{-4}$, integrated by parts & $0.076391 \pm 0.000006$ & $0.276422 \pm 0.000025$ \\
    Python (vegas), mean of three runs of $10^7$ points & $0.0761 \pm 0.0005$ & $0.2769 \pm 0.0011$ \\
    \bottomrule
  \end{tabular}
  \caption{FE 5D structure functions. Both Vegas results use the same
    integration variables (relative coordinates, see above) and $10^6$ points
    per iteration for each integral; they differ only in the integration by
    parts. Vegas errors are statistical. The Cuhre
    error is its own estimate, which covered the change to a ten times smaller
    $\epsilon_{\text{rel}}$. The Python error is the scatter of the three runs,
    about twice the error vegas reports for one run.}
  \label{tab:numbers}
\end{table}

\section{Implementation}

In the C++ code (\texttt{cpp/}), the boundary term is \texttt{FE4DIntegrand}
constructed with $Y = \ln x_0$, and the correction is
\texttt{FE5DCorrectionIntegrand} in \texttt{src/integrands.hpp}. They are
combined in \texttt{fe5d\_structure\_functions} in
\texttt{src/cross\_sections.hpp}. The option \texttt{-{}-no-by-parts} integrates
the original form \eqref{eq:sigma5d} instead.

The derivative is a central finite difference,
\begin{equation}
  \frac{\dd T}{\dd Y} \approx
  \frac{T(Y_+) - T(Y_-)}{Y_+ - Y_-} ,
  \qquad
  Y_\pm = \operatorname{clamp}\!\left(Y \pm h, \; Y_{\min}, \; Y_{\max}\right) ,
  \quad h = 10^{-4} ,
\end{equation}
taken as zero when $Y_+ \le Y_-$. It vanishes identically where the dipole is
frozen, and it needs nothing from the dipole model except its values and its
rapidity range, so it works with any interpolation of the BK solution. Its
error is about $10^{-8}$ relative.

A unit test (\texttt{tests/test\_sampling.cpp}) checks \eqref{eq:ibp} directly.
For fixed $u$, $u'$, $\theta$, $z$, it computes the original $\Mqq$ integral and
the boundary term plus the correction integral with high-order composite
Gauss--Legendre quadrature in $k$. The two agree to $10^{-7}$ relative, for both
polarizations and for finite and large $N_c$.

\end{document}
